The seventh phase of the Coupled Model Intercomparison Project (CMIP7),
the world's major contemporary project for running and comparing Earth System Models (ESMs),
is underway \citep{dunne2025evolving}.
A key part of model intercomparison is standardisation:
ensuring that assumptions which should be common across models are indeed common
so that the differences which emerge are isolated to factors of interest.
To support this standardisation,
a significant effort is invested in standardisation of model boundary conditions
and the importance of this standardisation is increasingly being recognised [TODO cite vaishali's paper].
These boundary conditions are referred to as forcings,
because they drive the model's response
(and play a role analogous to the role that force plays in Newton's second law
i.e. if one thinks e.g. in terms of a force and response model).

[TODO see if I can find some papers to refer to which give more context on how this process has been handled over time with a single sentence making that point e.g.
The process of generating forcings has evolved over multiple phases of CMIP.
There is now an explicit task team dedicated to the production, standardisation and documentation of these forcings datasets,
including papers like this one.]

For the world's most comprehensive climate models,
there are a large range of input forcings
[TODO cite various other forcing papers after "e.g.", also see if I can find a paper with a good overview of model inputs over time].
In this paper,
we document the atmospheric greenhouse gas concentration forcings to be used for the CMIP7 pre-industrial control (piControl) and historical experiments.
These forcings will also be used in a number of other experiments that we do not enumerate here.
Please see the papers describing these other CMIP7 experiments or
\newline \mbox{\url{https://wcrp-cmip.github.io/cmip7-guidance/docs/CMIP7/Experiment_set_up_and_Forcings/}}
for further details.
